\section{Implemented Surface}

The RCOUNT implementation adds audit-round and sample-draw records to
\texttt{rcount-core}. The package verifier now runs
\texttt{verify\_rla\_sampler\_replay} after CVR reconciliation, so a package
can only pass when the declared RLA sample draws replay.

The second implemented layer adds sample observations and a conservative rule:
any mark mismatch forces escalation. Every draw must have one observation. The
observation must identify the sampled CVR id, and its observed selections and
residual flags are compared to the normalized CVR row. If the number of
discrepancies is no greater than \texttt{max\_discrepancies}, the computed
status is \texttt{pass}; otherwise it is \texttt{escalate}.

Discrepancies are now first-class records. The verifier classifies mismatches
as selection, residual, combined selection-and-residual, or wrong-CVR
observations. When a package declares a discrepancy list, the declared
draw ids, CVR ids, and kinds must match the computed taxonomy exactly.

Margin metadata is also bound to the public summary before statistical stopping
math is attempted. The audit record names a winner selection, loser selection,
reported votes, reported margin, and diluted-margin denominator. The verifier
checks those fields against the canvassed jurisdiction total.

The first named statistical method is
\texttt{comparison-margin-threshold-v1}. It is intentionally conservative and
simple: it recomputes a risk estimate from sample size and reported margin,
adds a large penalty for declared discrepancies, checks the declared risk
estimate, and only permits \texttt{pass} when the estimate is at or below the
audit's risk limit and discrepancy count is within the declared threshold.

The first jurisdiction adapter is
\texttt{colorado-rule-25-comparison-v1}. It does not claim full Colorado
certification compliance. It verifies the RCOUNT-side contract needed by a
Colorado-style comparison audit fixture: a 20 digit public seed, the named
SHA-256 sampler, margin metadata, a comparison-margin stopping method, declared
risk estimate, and declared pass/escalate status.

The second jurisdiction adapter is \texttt{california-public-rla-v1}. It
targets a different public-contract layer: the audit record must name a
California-style ballot manifest format, public audit software identity, public
source URL, margin metadata, and declared audit status. This models the
open-algorithm and manifest-shape requirements before any later statutory risk
calculation is added.

\begin{center}
\small
\begin{tabularx}{\textwidth}{lX}
\toprule
Check & Failure caught \\
\midrule
\texttt{rla\_sampler\_replay} & Duplicate audit ids, invalid risk limits, empty or wrong-sized samples, unsupported algorithm ids, manifest hash drift, and edited sample draws. \\
\texttt{rla\_margin\_metadata} & Winner/loser selection ids, reported votes, reported margin, and diluted-margin denominator drift from the jurisdiction summary. \\
\texttt{rla\_stopping\_rule} & Missing or duplicate observations, wrong sampled CVRs, declared discrepancy taxonomy drift, risk-estimate drift, and declared stop/escalate outcomes that disagree with observed discrepancies. \\
\texttt{rla\_jurisdiction\_adapter} & Unsupported jurisdiction methods; Colorado-style records missing seed, sampler, margin, or stopping-risk fields; California-style records missing public software or manifest-format metadata. \\
\texttt{cvr\_summary\_total} & CVR rows still reconcile to the public summaries before they are used as the audit population. \\
\texttt{source\_hash\_match} & Source evidence bytes still match the declared source index. \\
\bottomrule
\end{tabularx}
\end{center}

The audit transcript maps sampler failures to
\texttt{rla\_sampler\_replay} and stopping failures to
\texttt{rla\_stopping\_rule}. Jurisdiction adapter failures surface as
\texttt{rla\_jurisdiction\_adapter}, making all three layers visible through
the same CLI surface as prior RCOUNT checks.
